\subsection{Quantum Measurements}
\label{sec:quantum-measurements}

A quantum measurement is the formal device by which one extracts
classical information from a quantum state. Each measurement is
specified by a finite set of \emph{outcomes}
$\{a\}_{a = 1}^{n}$ together with an operator on the Hilbert
space~$\mathcal{H}$ for each outcome; the probability of each outcome
is then read off from the state via the Born rule below. We
distinguish two levels of generality, projective measurements and the
strictly larger class of POVMs.

\begin{definition}[Projective measurement (PVM)]
\label{def:pvm}
A \emph{projective measurement} on a complex Hilbert space $\mathcal{H}$
of dimension $d$ is a finite family $\{\Pi_a\}_{a = 1}^{n}$ of
orthogonal projectors on $\mathcal{H}$,
\begin{equation}
  \Pi_a^{\dagger} \;=\; \Pi_a, \qquad
  \Pi_a \Pi_b \;=\; \delta_{ab}\, \Pi_a, \qquad
  \sum_{a = 1}^{n} \Pi_a \;=\; I.
  \label{eq:pvm}
\end{equation}
By the spectral theorem, projective measurements are in one-to-one
correspondence with self-adjoint observables
$A = \sum_a \lambda_a\, \Pi_a$ on~$\mathcal{H}$, the $\lambda_a$
being the eigenvalues attached to each outcome.
\end{definition}

\begin{definition}[Positive operator-valued measure (POVM)]
\label{def:povm}
A \emph{positive operator-valued measure} (POVM) on a complex Hilbert
space $\mathcal{H}$ with outcome set $\{1, \dots, n\}$ is a family
$\{M_a\}_{a = 1}^{n}$ of Hermitian operators satisfying
\begin{equation}
  M_a \;\succeq\; 0 \quad \forall\, a, \qquad
  \sum_{a = 1}^{n} M_a \;=\; I.
  \label{eq:povm-defining}
\end{equation}
Every projective measurement is a POVM (the $\Pi_a$ are positive,
Hermitian, and sum to the identity), but the converse fails: a POVM
element $M_a$ may have rank greater than one and the $M_a$ need not
mutually commute. POVMs are the most general notion of measurement
allowed by quantum mechanics; by Naimark's theorem every POVM arises
as a projective measurement on an enlarged Hilbert space, with the
extra degrees of freedom subsequently traced out.
\end{definition}

\begin{definition}[Born rule]
\label{def:born}
If the system is prepared in a state described by the density operator
$\rho$ and a POVM $\{M_a\}$ is performed, the probability of
observing outcome~$a$ is
\begin{equation}
  P(a \mid \rho) \;=\; \text{Tr}(\rho\, M_a).
  \label{eq:born}
\end{equation}
Positivity $M_a \succeq 0$ together with the resolution of the
identity $\sum_a M_a = I$ guarantees that
$\{P(a \mid \rho)\}_{a = 1}^{n}$ is a valid probability distribution
on the outcome set for every density operator~$\rho$.
\end{definition}

\paragraph{Indexed families of POVMs.}
In what follows we shall study not a single POVM but \emph{collections}
of POVMs sharing the same Hilbert space and outcome set, indexed by a
classical measurement \emph{setting}~$x \in \{1, \dots, m\}$. We
adopt the standard notation $M^{(x)} = \{M_{a|x}\}_{a = 1}^{n}$ for
the POVM associated with setting~$x$, so that each $M_{a|x}$ is a
Hermitian operator on $\mathcal{H}$ satisfying
\eqref{eq:povm-defining} for every fixed~$x$, and the Born rule
extends naturally to $P(a \mid x, \rho) = \text{Tr}(\rho\, M_{a|x})$.
The central question of this report---developed in the next
subsection and formalised as a semidefinite program in
Section~\ref{sec:sdp-2qubit}---is to what extent such a collection
$\{M^{(x)}\}_{x}$ can be realised \emph{simultaneously} by a single
physical measurement.

\begin{definition}[White-noise mixture of a POVM]
\label{def:white-noise}
Given a POVM $\{M_a\}$ on $\mathcal{H} = \mathbb{C}^{d}$ and a noise
parameter $\eta \in [0, 1]$, the corresponding \emph{$\eta$-noisy}
POVM has elements
\begin{equation}
  M^{\eta}_{a} \;=\; \eta\, M_{a} \;+\; (1 - \eta)\, \text{Tr}(M_a)\,
                     \frac{I}{d}.
  \label{eq:noisy-povm}
\end{equation}
At $\eta = 1$ the original sharp measurement is recovered; at
$\eta = 0$ every outcome operator becomes proportional to the
maximally mixed state $I/d$, so the outcome distribution
$P(a \mid \rho)$ no longer depends on $\rho$ and the measurement is
uninformative. The parameter $\eta$ will serve as the quantitative
figure of merit of measurement incompatibility throughout the rest of
the report.
\end{definition}
